\section{Introduction}

Reconstructed driving scenes support sensor simulation, perception analysis, and counterfactual replay.
Recent neural simulators represent complex appearance and moving objects with impressive fidelity
\cite{yang2023unisim,zheng2024lidar4d,wu2024dynfl}. Physical queries impose an additional requirement:
the reconstructed surface must agree with where a sensor first encounters an object.
A visually plausible vehicle can still extend into measured free space, while a sparse but accurate point set
can miss much of its surface. These errors affect different parts of the same ray and call for joint surface
and return evaluation.

Two challenges arise in learning such a representation. First, multi-frame observations of moving Actors must
be aligned before they become completion targets. A small collection of seeds also has insufficient capacity
to cover a dense target surface. Second, a primitive's location and its contribution to a sensor return are
different variables. The same anchor may receive occupied evidence in one view and free-space evidence in another.
Assigning every overlapping Gaussian unit opacity makes the return depend on sampling density, while a single
confidence score discards the structure of these observations.

We address these challenges with \emph{Evidential Actor Surfaces} (EAS), illustrated in
Fig.~\ref{fig:eas-overview}. Each tracked Actor has a canonical physical surface, continuous return evidence,
and a separately owned rigid trajectory. We construct build inputs and disjoint supervision targets by
transforming LiDAR sweeps into the Actor frame. A set decoder expands every completion seed into four child
surfaces. Set distance, local-plane and scale targets, frame coverage, and ray losses jointly supervise the
children while observed anchors stay fixed. This provides a direct route from measured geometry to the deployed surface.

For sensor termination, we retain the observation history at the corpus producer and learn separate evidence
heads for anchors and children. Their free, occupied, and unknown masses summarize multi-view support and
contradiction. On the fixed surface, occupied mass weights a normalized categorical distribution over ray depth.
This composition yields a useful first-return measure from the learned evidence. Its normalization also permits
an exact component-wise analysis: decreasing a component's weight moves the pre-hit probability toward the
remaining mixture. We derive the finite change, identifying the condition under which attenuation decreases
early-return probability.

The dynamic scene interface follows the same separation of variables. A read-only SE(3) pose places the canonical
surface in the world; a sibling visual state carries appearance. This makes rigid composition equivariant and
prevents image optimization from moving the physical surface. The representation therefore combines learnable
local geometry with explicit control over which state each observation may update.

Our experiments separate geometry learning, return composition, and coordinate ownership. On the nuScenes
development fold, surface supervision improves hazardous early returns, hit recall, and Chamfer jointly.
Categorical evidence improves early returns and hit recall across the all, hazardous, and clear strata on
the same geometry. A frozen AV2 cohort measures cross-sensor behavior, and analytical and numerical audits
connect the return measure to its attenuation response. We make three contributions:
\begin{itemize}
\item an Actor-canonical completion representation trained with target-set, frame, and ray supervision;
\item a producer-evidential categorical return measure with an exact finite attenuation identity;
\item a factorized physical--pose--visual interface, evaluated through surface accuracy, first-return behavior,
and SE(3) composition.
\end{itemize}
